\section{The theory of Ng et al.}
\label{sec:ng}

We restate what is needed from \cite{ng2025}. For each environment $u\in[m]$ there is a latent vector
$Z^u=(Z_1^u,\dots,Z_n^u)$ and an observation $X^u=g(Z^u)$, with $g$ a $\mathrm{C}^2$ 
\textbf{\textit{diffeomorphism onto its image ((injective, smooth, smooth inverse)}}, shared by all environments. Each coordinate satisfies
$Z_i^u=f_i^{(u)}(\PA(Z_i^u;\G_Z))+\epsilon_i^{(u)}$ with $\epsilon_i^u$ Gaussian (additive noise
model, ANM; the meaning of $\epsilon_i^{(u)}$ is discussed in \cref{app:eps}), where the same unknown
DAG $\G_Z$ holds in every environment and the mechanisms $f_i^{(u)}$ may change arbitrarily with
$u$. Write $p^{(u)}_Z$ for the density of $Z^u$; it is assumed third-order differentiable (and hence 
--  continuous) with \emph{full support on $\R^{n}$}.

\subsection{Ng et al. objective}
Ng et al. objective is to learn an \textit{observationally equivalent} model 
$(\hat g, p_{\hat Z}, \hat \G_{\hat Z})$
that follows the above data generating process, and generates a data distribution that matches the 
distribution of the observed random variables $X$ across the given environments, i.e. 
(\cite[Equation 2]{ng2025}), 
\begin{equation}
p^{(u)}_{\hat X}(x) = p^{(u)}_X(x),\ \forall u \in [m],\ x \in X^{(u)},
\label{eq:obs_equiv}
\end{equation}
where $\hat Z$ denote the estimated latent variables and $\hat X = \hat{g}(\hat Z)$. In this case, 
the models $(\hat g, p_{\hat Z}, \G_{\hat Z})$ and $(g, p_Z, \G_Z)$ are said to be 
\textit{observationally equivalent}.

\subsection{Graph notions}
The \emph{Markov network} $\Mk_Z$ has an edge $\{Z_i,Z_j\}$ iff
$Z_i\not\perp Z_j\mid Z_{[n]\setminus\{i,j\}}$, equivalently iff \linebreak
$\partial^2\log p^{(u)}(Z)/\partial Z_i\partial Z_j\not\equiv0$ for some $u$; under faithfulness it is
the moral graph of $\G_Z$. Write $\mathcal{E}(\Mk_Z)$ for its edge set, $\Delta(\Mk_Z)$ for its 
set of triangles and $\deg_{\Mk_Z}(v)$ for the degree of a node. 
The \emph{surrounding parents} of $Z_i$ are
\[
\sur(Z_i;\G_Z):=\{Z_j\in\PA(Z_i;\G_Z)\mid Z_j\text{ is a parent of every child of }Z_i\}.
\]

\subsection{Sufficient-change vectors}
For a set $Z'$ of latents that contains all of its own ancestors (an \emph{ancestral} set), with
$n'=|Z'|$, let $\G_{Z'}$ be the induced DAG, $\Mk_{Z'}$ its Markov network, and define
\begin{align}
\vect\tau(Z',u) &= \Bigl(\tfrac{\partial\log p^{(u)}(Z')}{\partial Z_i}\Bigr)_{Z_i\in Z'}
\oplus\Bigl(\tfrac{\partial^2\log p^{(u)}(Z')}{\partial Z_i^2}\Bigr)_{Z_i\in Z'}
\oplus\Bigl(\tfrac{\partial^2\log p^{(u)}(Z')}{\partial Z_i\partial Z_j}\Bigr)_{\{Z_i,Z_j\}\in \mathcal{E}(\Mk_{Z'})},
\label{eq:tau}\\[2pt]
\vect w(Z',u) &= \vect\tau(Z',u)
\oplus\Bigl(\tfrac{\partial^3\log p^{(u)}(Z')}{\partial Z_i^3}\Bigr)_{Z_i\in Z'\setminus\Sk(\G_{Z'})}
\oplus\Bigl(\tfrac{\partial^3\log p^{(u)}(Z')}{\partial Z_i^2\partial Z_j}\Bigr)_{\substack{\{Z_i,Z_j\}\in \mathcal{E}(\Mk_{Z'}),\\ Z_i\notin\Sk(\G_{Z'})}} \nonumber \\
& \quad \oplus\Bigl(\tfrac{\partial^3\log p^{(u)}(Z')}{\partial Z_i\partial Z_j\partial Z_k}\Bigr)_{\{Z_i,Z_j,Z_k\}\in\Delta(\Mk_{Z'})}.
\label{eq:w}
\end{align}
In the fourth group in \cref{eq:w} we read the index set as \emph{ordered} pairs $(i,j)$: an edge contributes one entry for each endpoint that is not a sink, taking that endpoint as the squared variable. This is how
the entries arise in Eq.~(7) in the proof of Proposition~2 of \cite{ng2025}, where the sum runs over
such pairs, and how the term $3|E|$ in their count arises. (The fourth group is printed in the paper with
the index condition ``$i<j$''; see \cref{rem:ij}.)

\begin{assumption}[Sufficient changes for the Markov network, \cite{ng2025}]\label{as:1}
For every ancestral $Z'$ and every value of $Z'$ there exist $2n'+|\mathcal{E}(\Mk_{Z'})|+1$ values
$u_0,u_1,\dots$ such that the vectors $\vect\tau(Z',u_j)-\vect\tau(Z',u_0)$, $j\ge1$, are linearly
independent.
\end{assumption}

\begin{assumption}[Sufficient changes for sink nodes, \cite{ng2025}]\label{as:2}
For every ancestral $Z'$ and every value of $Z'$, with $L:=|\vect w(Z',u)|$, there exist $L+1$ values
$u_0,\dots,u_L$ such that the vectors $\vect w(Z',u_j)-\vect w(Z',u_0)$, $j=1,\dots,L$, are linearly
independent.
\end{assumption}

\begin{theorem}[Identifiability with latent ANMs, \cite{ng2025}]\label{thm:ng}
Under the data generating process above, \cref{as:1,as:2} and faithfulness, Algorithm~1 of 
\cite{ng2025}, which is required to find an \textit{observationally equivalent} model of the data
generating process, returns $(\hat{\G},\hat Z)$ such that, for some permutation $\pi$ of the estimated 
latents, (i) $\hat{\G}_\pi$ and $\G_Z$ are identical, and 
(ii) each $\hat Z_{\pi(i)}$ is solely a function of a subset of
$\{Z_i\}\cup\sur(Z_i;\G_Z)$.
\end{theorem}

The role of the third derivatives is to identify sink nodes: for a sink $Z_i$ with \textit{Gaussian}
noise, $\partial^3\log p/\partial Z_i^2\partial Z_j=0$ for all $j$ \cite{ng2025} (a property used in
score matching for observed variables in \cite{rolland2022}). With the Markov network known, the parents
of an identified sink are its neighbors; the sink is frozen and the procedure repeats on the rest.
Without further assumptions nonlinear ICA and hence CRL are not identifiable \cite{hyvarinen1999};
the environments supply the needed variation. The Markov-network stage alone is the result of
\cite{zhang2024}, which yields identification only up to mixing with ``intimate neighbors''.

Algorithm~1 requires learning $(\hat g^t, p^t_{\hat Z}, \G^t_{\hat Z})$ where $t$ is the iteration 
index under additional constraints to achieve \cref{eq:obs_equiv} -- observational equivalence. 
They propose to use a VAE to that end with additional regularization terms. 

\subsection{Gaussianity assumptions in ANM}
Ng et al. assume Gaussianity in two places: (i) In their theory exogenous noise terms and (ii) in the implementation of the VAE where all distributions that need to be specified (encoder, decoder, prior) are assumed Gaussian.

\begin{itemize}
\item \textit{In the theory}: Gaussianity of the exogenous noise is used for exactly one purpose -- the 
sink-detection property. For a sink $Z_i$ under Gaussian noise, $\log p(Z_i\mid\text{parents})$ is 
quadratic in $Z_i$, so $\partial^3\log p/\partial Z_i^2\partial Z_j=0$ for all $j$ (ANM case). This 
vanishing is what lets Algorithm~1 identify sinks via third-derivative tests and peel them off one at a 
time. It's genuinely Gaussian-specific (quadratic log-density $\Leftrightarrow$ Gaussian) -- no other 
noise family gives this cleanly. Nothing else in Assumptions~1/2 (the linear-independence conditions) 
needs Gaussianity per se; they're stated generically in terms of derivatives of $\log p^{(u)}$.

\item \textit{In the implementation}: Only one of the three Gaussian choices actually follows from the 
theory:
\begin{itemize}
\item \textbf{Prior over $\hat Z$} (\texttt{mean/logvar} nets, ANM structure) -- follows from the 
theory. This is the direct implementation of the Gaussian exogenous-noise assumption, and it's what the 
sink-freezing logic is trying to exploit (even if, per \cref{app:code}, the actual freezing rule is a 
sparsity-mask heuristic, not the literal third-derivative test).
\item \textbf{Encoder $q(z\mid x)$ Gaussian} -- does not follow from the theory. The paper says nothing 
about the recognition model's distribution; this is a standard VAE/reparameterization convenience, 
orthogonal to the noise-model assumption.
\item \textbf{Decoder Gaussian $\rightarrow$ MSE reconstruction loss} -- does not follow from the theory 
either. The paper never states a reconstruction loss at all (already noted in \cref{app:code}) as
``\textit{an implicit modelling choice not discussed in the paper}''.
This is purely a VAE-training-objective choice, not an implication of Gaussian noise in the SEM.
\end{itemize}
So: theory constrains the prior; encoder and decoder Gaussianity are independent engineering choices 
riding along with ``\textit{it's a VAE}'', not consequences of the identifiability assumption.
\end{itemize}

\subsection{Gaussianity assumptions in HNM}
\begin{itemize}
\item \textit{Theory}: Gaussianity still does exactly the same job -- identifying sinks -- but the 
property it buys is weaker. Now $\epsilon_i\sim\mathcal N(0,1)$ and the scale $\sigma_i^{(u)}(\PA(Z_i))$ 
carries the parent-dependence, so $\log p(Z_i\mid\PA)$ is no longer purely quadratic in $Z_i$ 
(the $\log\sigma_i(\PA)$ term depends on the parents too). Result: $\partial^3\log p/\partial 
Z_i^2\partial Z_j=0$ only for $Z_j\notin\PA(Z_i)$ -- vanishes for non-parents, not for parents. So 
Gaussianity of $\epsilon_i$ is still what makes the non-parent third derivatives vanish; it just no 
longer kills the parent ones, because the heteroscedastic scale reintroduces parent-dependence into the 
third derivative.

Practically this means the sink test still works, but is less informative: you learn 
\textit{``no non-parent neighbors act on $Z_i$ nontrivially,}'' but you don't get the same clean 
``\textit{this column of the adjacency is zero}'' signal from the pure log-density curvature -- you 
still need the Markov network (second derivatives) to know who the candidate neighbors are, and the 
third-derivative sink test now only rules out non-parents rather than certifying the full parent set 
directly the way it did in the ANM.
\newpage
\item \textit{Implementation}:
\begin{itemize}
\item \textbf{Prior} -- same story, just needs \texttt{log\_var\_net} (or whatever computes 
$\hat\sigma_i^{(u)}$) to also take the masked/estimated parents as input, not just $u$. Still 
{theory-driven}.
\item \textbf{Encoder/decoder Gaussianity} -- unaffected, same VAE convenience as before, orthogonal to 
ANM vs. HNM.
\item \textbf{Sink-freezing heuristic} (\texttt{find\_sinknodes\_and\_fix}) -- 
this one actually gets harder to 
justify under HNM, since the mask-based heuristic was already a loose proxy for the ANM's sharp third-
derivative test; under HNM the thing it's supposed to approximate is itself weaker, so the gap between 
``\textit{heuristic zeros out a column}'' and ``\textit{theorem's sink property}'' widens rather than 
narrows.
\end{itemize}
\end{itemize}